\scanpage{9}

Formes \uline{quadratiques binaires et} revêtements quadratiques

1. Forme quadratique à scalaire près, structure de similitude.

1.1. $(X,\underline{\mathcal O}_X)$ désigne un schéma, plus généralement un topos loc. annelé. Soit $E$ un Module sur $X$,

\[\underline{\mathrm{Quad}}(E)\simeq\underline{\mathrm{Hom}}(\Gamma^2(E),E)\]

le faisceau des formes quadratiques sur $E$. Si $E$ est loc. libre de type fini, on a un iso. can.

\[\underline{\mathrm{Quad}}(E)\simeq\underline{\mathrm{Sym}}^2(E^\vee),\]

(en général on a seulement un homomorphisme can.

\[\underline{\mathrm{Sym}}^2(E^\vee)\longrightarrow\underline{\mathrm{Quad}}(E),\]

qui transforme $f.g$ ($f,g\in\Gamma E^\vee$) dans $Q_{\unsure{fg|f,g}}:x\longmapsto f(x)g(x)$ sur $E$ \ldots).

Le faisceau \deleted{$\unsure{\underline{\mathcal O}}$} en groupes $\underline{\mathcal O}_X^\ast$ opère sur $\underline{\mathrm{Quad}}(E)$, d'où un faisceau quotient

\[\underline{\mathrm{Sim}}(E)\simeq\underline{\mathrm{Quad}}(E)/\underline{\mathcal O}^{\ast}.\]

Une section du faisceau quotient s'appelle une \uline{structure de similitude} sur $E$. Une telle section correspond à un sous-faisceau ensembliste de $\underline{\mathrm{Quad}}(E)$, qui \deleted{loc} soit à fibres non vides et sur lequel $\mathcal O^\ast$ opère transitivement. Localement, une telle structure est définie par une section $Q$ de $\underline{\mathrm{Quad}}(E)$ i.e. par une forme quadratique sur $E$, \deleted{\unsure{deux}} deux telles sections $Q$, $Q'$ définissent la $=$ structure de similitude \unsure{sss|ssi} $\exists$ localement $a\in\Gamma\underline{\mathcal O}_X^\ast$, avec $Q'=aQ$.

Le cas le plus intéressant est celui où \deleted{$a\mapsto aQ$}

\[a\longrightarrow a.Q:\underline{\mathcal O}_X\ (\text{sans $\ast$\,!})\longrightarrow\underline{\mathrm{Quad}}(E)\]

est injectif. Dans ce cas, \deleted{si} $Q'$ \deleted{\unsure{est connu de|définit la même}} \inserted{définit la} $=$
% archive sheet number: 8
